\section{Рабочий проект}
\subsection{Классы, используемые при разработке программы}

Можно выделить следующий список классов и их методов, использованных при разработке программно-информационной системы (таблица \ref{class:table}).

\renewcommand{\arraystretch}{0.8} % уменьшение расстояний до сетки таблицы
\begin{xltabular}{\textwidth}{|X|p{2.5cm}|>{\setlength{\baselineskip}{0.7\baselineskip}}p{4.85cm}|>{\setlength{\baselineskip}{0.7\baselineskip}}p{4.85cm}|}
\caption{Описание классов, используемых в программе\label{class:table}}\\
\hline \centrow \setlength{\baselineskip}{0.7\baselineskip} Название класса & \centrow \setlength{\baselineskip}{0.7\baselineskip} Модуль, к которому относится класс & \centrow Описание класса & \centrow Методы \\
\hline \centrow 1 & \centrow 2 & \centrow 3 & \centrow 4\\ \hline
\endfirsthead
\caption*{Продолжение таблицы \ref{class:table}}\\
\hline \centrow 1 & \centrow 2 & \centrow 3 & \centrow 4\\ \hline
\finishhead
IN & Модуль авторизации & IN –  класс программы, проводящий аутентификацию пользователя в системе. После ввода кода сотрудника, данные проверяются на корректность. При успешной авторизации, работник становится текущим пользователем. В противном случае, возникает сообщение об ошибке.
& add\_symbol(digit) - метод добавления символа в строку. 
check\_password() - проверяет пароль на корректность, из базы данных импортирует соответствующую должность.
open\_manager\_window
() - сравнивает должность из таблицы БД с требуемой, при успехе инициализируется класс ManagerWindow.
open\_shef\_window() - сравнивает должность из таблицы БД с требуемой, при успехе инициализируется класс ManagerMenu.
open\_off\_window() - сравнивает должность из таблицы БД с требуемой, при успехе инициализируется класс Tabel.\\
\hline ManagerWin-
dow & Главный модуль для управляющего & ManagerWindow – Класс для создания окна, состоящего из кнопок для работы с модулями управления сотрудниками, столами, скидками, типами оплаты. & button\_manager() - создание кнопок для инициализации классов AddSale, AddTabels, AddPayType, ManagerStaff, ManagerPost.
open\_window(index) открывает окно, содержащее выбранный класс.\\
\hline Manager-
Staff & Модуль управления сотрудниками & ManagerStaff – Класс для работы с данными сотрудников и должностей.& save\_data() - функция для сохранения данных в таблицу БД.
tabel(self) - создание таблицы в окне программы.
add\_data(self) - добавление данных в таблицу программы.
edit\_data(self) - редактирование данных в таблице программы.
delete\_data(self) - удаление данных из таблицы программы.
window\_search\_data
(self) - поиск данных в таблице программы.
refresh\_data(self) - обновление таблицы программы, данные импортируются из БД.\\

\hline AddPayTy-
pes & Модуль управления типами оплаты & AddPayTypes – Класс для работы с данными типов оплаты. & 
add\_data(self) - сохранение введенных данных в таблицу БД.
clear\_fields(self) - очищает поля, после сохранения данных.\\

\hline AddTabels & Модуль управления столами & AddTabels – Класс для работы с данными столов. & add\_data(self) - сохранение введенных данных в таблицу БД.
clear\_fields(self) - очищает поля, после сохранения данных.\\

\hline AddSale & Модуль управления скидками & AddSale – Класс для роботы с данными скидкок. & add\_data(self) - сохранение введенных данных в таблицу БД.
clear\_fields(self) - очищает поля, после сохранения данных.\\



%\hline CFile & Главный модуль & CFile – Класс для работы с файлами и изображениями & array GetFileArray (int file\_id)
%Метод возвращает массив, содержащий описание файла (путь к файлу, имя файла, размер) с идентификатором file\_id\\
\end{xltabular}



\renewcommand{\arraystretch}{1.0} % восстановление сетки

\subsection{Модульное тестирование разработанного web-сайта}

Модульный тест для класса User из модели данных представлен на рисунке \ref{unitUser:image}.

\begin{figure}[ht]
\begin{lstlisting}[language=Python]
from django.test import TestCase
from .models import *
User = get_user_model()


class ShpoTestCases(TestCase):

    def setUp(self) -> None:
        self.user = User.objects.create(username='testtestovich', password='testtestovich', first_name='Sad', last_name='')

    def test_2(self):

        self.assertEqual(self.user.first_name, 'Sad')
        self.assertEqual(self.user.last_name, 'Cat')
        print((self.user))
        print((self.user.first_name))
        print((self.user.last_name))
\end{lstlisting}  
\caption{Модульный тест класса User}
\label{unitUser:image}
\end{figure}

\subsection{Системное тестирование разработанного web-сайта}

На рисунке  представлена главная страница сайта «Русатом – Аддитивные технологии».
\newpage % при необходимости можно переносить рисунок на новую страницу


На рисунк представлен динамический вывод заголовков, включающий в себя искомые фразы при поиске фраз.


На рисунк представлен ввод данных для публикации новости.

